\subsection{Tangent Planes and Linear Approximations}
\content{

}{% presentation
\vspace{3in}
}{% nopresentation
}{% comments
Derive the equation of the tangent plane. This ends up being 
$$ z-z_0 = f_x(a,b)(x-a)+f_y(a,b)(y-b). $$
}

\content{
\begin{example} Find the tangent plane to the elliptic paraloloid 
$$ z = 2x^2+y^2. $$
\end{example} 
}{% presentation
\vspace{3in}
}{% nopresentation
}{% comments
}

\content{
The tangent plane can be used to approximate values of a function near a point
whose value is known:
\begin{example} Use the tangent plane to approximate $f(1.1,-0.1)$ where 
$$ f(x,y) = xe^{xy}. $$
\end{example}
}{% presentation
\vspace{3in}
}{% nopresentation
}{% comments
approx. $1.1e^{-0.11} \approx 0.98542$.
}

\subsubsection{Differentials}

\content{
Recall that for functions of a single variable, we have that the differential
$dy$ is given by 
$$ dy = f'(x)dx. $$
We have a similar definition for functions of two variables: 
\begin{definition} We define the differential $dz$, also called the total
differential, to be the value 
$$ dz = f_x(x,y)dx + f_y(x,y)dy=\frac{\partial f}{\partial x}dx + \frac{\partial f}{\partial y}dy. $$
\end{definition}
The Linearization of $f$ at $(a,b)$ can then be written 
$$ f(x,y) \approx f(a,b) + dz. $$
}{% presentation
}{% nopresentation
}{% comments
}

\content{
\begin{example} If 
$$ z = f(x,y) = x^2+3xy - y^2, $$
find the differential $dz$.
\end{example}
}{% presentation
\vspace{3in}
}{% nopresentation
}{% comments
}

\content{
\begin{example} The baes radius and height of a right circular cone are measured
as 10 cm and 25 cm, respectively, with a possible error in measurement of as
much as 0.01 cm in each. Use differentials to estimate the maximum error in the
calculated volume of the cone. 
\end{example}
}{% presentation
\vspace{4in}
}{% nopresentation
}{% comments
}
